\documentclass[12pt,a4paper]{article}
\usepackage[a4paper, margin=2.5cm]{geometry}
\usepackage[utf8]{inputenc}
\usepackage[T1]{fontenc}
\usepackage[magyar]{babel}
\usepackage{amsmath, amssymb, amsthm}
\usepackage{xcolor}
\usepackage{tcolorbox}
\tcbuselibrary{theorems,skins,breakable}
\usepackage{tikz}
\usepackage{enumitem}

%%% Megjelenítéshez automatikusan kiegészített csomagok %%%
\usepackage{pgfplots}
\pgfplotsset{compat=1.18}
\usepackage{graphicx}
\usepackage{booktabs}
\usepackage{float}
\usepackage{caption}
\usepackage{microtype}

\usepackage{hyperref}

\title{Permutáció}
\author{Matematika oktatási segédlet}
\date{}

\begin{document}

\maketitle

\section*{Elméleti áttekintés és képletek}

A kombinatorika egyik alapvető kérdése, hogy adott elemeket hányféleképpen tudunk sorba rendezni. Ezt a folyamatot permutációnak nevezzük.

\begin{tcolorbox}[colback=gray!5!white,colframe=gray!40!black,title=1. Ismétlés nélküli permutáció, fonttitle=\bfseries, arc=2mm, boxrule=0.5mm]
Ha $n$ darab \textbf{különböző} elemünk van, akkor ezek összes lehetséges sorrendjének (azaz ismétlés nélküli permutációjának) a száma ($P_n$):
\begin{equation*}
    P_n = n! = n \cdot (n - 1) \cdot (n - 2) \cdot \dots \cdot 1
\end{equation*}
Megállapodás szerint a nulla faktoriális értéke: $0! = 1$.
\end{tcolorbox}

\begin{tcolorbox}[colback=gray!5!white,colframe=gray!40!black,title=2. Ismétléses permutáció, fonttitle=\bfseries, arc=2mm, boxrule=0.5mm]
Ha $n$ darab elemünk van, de ezek között vannak azonosak is ($k_1$ darab az első fajtából, $k_2$ darab a második fajtából, $\dots$, $k_r$ darab az $r$-edik fajtából, ahol $k_1 + k_2 + \dots + k_r = n$), akkor az ismétléses permutáció száma:
\begin{equation}
    P_n^{k_1, k_2, \dots, k_r} = \frac{n!}{k_1! \cdot k_2! \cdot \dots \cdot k_r!}
\end{equation}
\end{tcolorbox}

\vspace{0.3cm}
\noindent\textbf{Fontos módszertani tippek a feladatokhoz:}
\begin{itemize}
    \item \textbf{Szomszédsági feltétel („ragasztás”):} Ha bizonyos elemeknek egymás mellett kell lenniük, gondolatban ragasszuk őket egyetlen egységgé (csomaggá), rendezzük a csomagokat, majd szorozzunk be a csomagon belüli elemek lehetséges helycseréivel ($k!$).
    \item \textbf{Kizárásos feladatok:} Ha azt vizsgáljuk, hogy valami \emph{nem} történhet meg, gyakran egyszerűbb az összes esetből kivonni a tiltott (kívánt ellentétes) esetek számát.
    \item \textbf{Nulla a számjegyek elején:} Határozottan figyeljünk arra, hogy egy természetes szám legelső helyi értékére (pl. ötjegyű szám tízezresek helyére) a $0$ nem kerülhet!
\end{itemize}

\newpage

\section*{Gyakorló feladatok}

\begin{enumerate}[label=\arabic*., leftmargin=*]
    
    \item \textbf{1. Feladat} \textsf{\small [K1 – Alap]} \\
    Egy öttagú baráti társaság (Anna, Bea, Cecil, Dani, Efraim) megérkezik a színházba, és egy öt egymás melletti székből álló sorba ül le.
    \begin{enumerate}[label=\alph*)]
        \item Hányféleképpen ülhetnek le egymás mellé?
        \item Hány olyan ülésrend létezik, amelyben Anna ül a bal szélső (1.) széken?
    \end{enumerate}

    \vspace{0.4cm}

    \item \textbf{2. Feladat} \textsf{\small [K1 – Alap]} \\
    Határozd meg az alábbi kifejezések pontos értékét számológép nélkül:
    \begin{enumerate}[label=\alph*)]
        \item $\frac{8!}{6!}$
        \item $\frac{10!}{8! \cdot 2!}$
        \item $\frac{7! - 6!}{5!}$
    \end{enumerate}

    \vspace{0.4cm}

    \item \textbf{3. Feladat} \textsf{\small [K1 – Alap]} \\
    Az \textbf{ANANÁSZ} szó betűit sorba rendezzük. 
    \begin{enumerate}[label=\alph*)]
        \item Hány betűs a szó, és mennyi az egyes betűk gyakorisága (leltár)?
        \item Hány különböző (értelmes vagy értelmetlen) szót tudunk kirakni a betűiből?
    \end{enumerate}

    \vspace{0.4cm}

    \item \textbf{4. Feladat} \textsf{\small [K2 – Közepes]} \\
    Egy könyvespolcon 6 különböző regény sorakozik egymás mellett.
    \begin{enumerate}[label=\alph*)]
        \item Hányféleképpen helyezhetjük el őket a polcon?
        \item Hány olyan elrendezés létezik, amelyben két adott regény (az A és a B) pontosan egymás mellett van?
        \item Hány olyan elrendezés van, amelyben az A és a B regény \emph{nem} kerül egymás mellé?
    \end{enumerate}

    \vspace{0.4cm}

    \item \textbf{5. Feladat} \textsf{\small [K2 – Közepes]} \\
    Írjuk fel az összes olyan négyjegyű számot, amely a $1, 2, 3, 4$ számjegyek felhasználásával készül (minden számjegyet pontosan egyszer használunk fel). 
    \begin{enumerate}[label=\alph*)]
        \item Hány darab ilyen négyjegyű szám létezik összesen?
        \item Hány olyan szám van ezek között, amely páros?
        \item Hány olyan szám van, amely 2500-nál nagyobb?
    \end{enumerate}

    \vspace{0.4cm}

    \item \textbf{6. Feladat} \textsf{\small [K2 – Közepes]} \\
    Egy büfében 4 fiú és 4 lány áll sorban a pénztárnál.
    \begin{enumerate}[label=\alph*)]
        \item Hányféleképpen állhatnak sorba, ha nincs semmilyen megkötés?
        \item Hányféleképpen állhatnak sorba, ha az összes fiú egymás mellett, és az összes lány is egymás mellett áll?
        \item Hány olyan sorrend létezik, amelyben a fiúk és a lányok felváltva állnak (azaz két azonos nemű személy nem áll közvetlenül egymás mögött)?
    \end{enumerate}

    \vspace{0.4cm}

    \item \textbf{7. Feladat} \textsf{\small [K2 – Közepes]} \\
    A $0, 1, 1, 2, 2, 2$ számjegyek felhasználásával hatjegyű számokat képezünk, mind a 6 számjegyet felhasználva egy-egy számban.
    \begin{enumerate}[label=\alph*)]
        \item Hány darab hatjegyű szám készíthető összesen?
        \item Hány olyan szám van a fenti adatokból, amely nem kezdődhet nullával (azaz valódi hatjegyű szám)?
    \end{enumerate}

    \vspace{0.4cm}

    \item \textbf{8. Feladat} \textsf{\small [K3 – Nehezebb]} \\
    Hat barát (három házaspár: Anna és Péter, Eszter és Gábor, Dóra és Tamás) kerek asztal köré ül vacsorázni. A forgatással egymásba átvihető ülésrendek azonosnak számítanak.
    \begin{enumerate}[label=\alph*)]
        \item Hányféleképpen ülhetnek le a kerek asztal köré összesen?
        \item Hány olyan ülésrend létezik, amelyben minden férj a felesége mellett ül?
    \end{enumerate}

    \vspace{0.4cm}

    \item \textbf{9. Feladat} \textsf{\small [K3 – Nehezebb]} \\
    Egy sakktornán 5 magyar és 3 külföldi versenyző indul. A díjkiosztón egyetlen hosszú sorban állnak fel a fotózáshoz.
    \begin{enumerate}[label=\alph*)]
        \item Hányféle sorrendben állhatnak fel, ha a magyarok és a külföldiek is tetszőleges sorrendben követhetik egymást?
        \item Hány olyan felállási sorrend létezik, amelyben a szélsőséges helyeken (a sor két végén, azaz az 1. és a 8. helyen) egyaránt magyar versenyző áll?
    \end{enumerate}

    \vspace{0.4cm}

    \item \textbf{10. Feladat} \textsf{\small [K3 – Nehezebb]} \\
    Adott az $\{1, 2, 3, 4, 5\}$ számhalmaz. Ebből a számjegyekből tetszőleges sorrendben ötjegyű számokat képezünk úgy, hogy mind az 5 számjegyet pontosan egyszer használjuk fel.
    \begin{enumerate}[label=\alph*)]
        \item Mennyi az összes így képezhető ötjegyű szám összege? (\emph{Útmutatás: Gondold végig, hogy az egyes helyi értékeken – egyesek, tízesek, stb. – az 1, 2, 3, 4, 5 számjegyek egyaránt hányszor fordulnak elő!})
    \end{enumerate}

\end{enumerate}

\end{document}